% Motor: XeLaTeX. Sesión 29 - material docente UTP.
% Estructura docente de la Sesión 28.
\documentclass[aspectratio=169,10pt,c]{beamer}
% Reutiliza sin cambios la identidad visual UTP validada en la Sesion 08.
% Reutiliza sin cambios la identidad visual UTP validada en la Sesion 05.
% La segunda ruta permite compartir los assets comunes sin duplicar binarios.
% !TEX root = slides.tex
% Preambulo comun: estilo UTP, colores, fuentes, plantillas de beamer y
% comandos de apoyo didactico. Se carga desde slides.tex tras \documentclass.

% La opción "c" de beamer reparte el espacio libre 1:1.5 (arriba:abajo), por lo
% que el contenido queda por encima del centro. Igualamos ambos rellenos.
\makeatletter
\define@key{beamerframe}{c}[true]{%
  \beamer@frametopskip=0pt plus 1fill\relax%
  \beamer@framebottomskip=0pt plus 1fill\relax%
  \beamer@frametopskipautobreak=0pt plus .5\paperheight\relax%
  \beamer@framebottomskipautobreak=0pt plus .5\paperheight\relax%
}
\makeatother

\usepackage{graphicx}
\usepackage{tikz}
\usepackage{xcolor}
\usepackage{iftex}

\ifPDFTeX
  \usepackage[utf8]{inputenc}
  \usepackage[T1]{fontenc}
  \usepackage{lmodern}          % tipografia escalable para el respaldo con pdflatex
  \renewcommand{\familydefault}{\sfdefault}
\else
  \usepackage{fontspec}
  \IfFontExistsTF{Arial}
    {\setsansfont{Arial}[BoldFont={Arial Bold}]}
    {\setsansfont{Helvetica}[BoldFont={Helvetica Bold}]}
\fi

\usepackage[spanish,es-nodecimaldot]{babel}
\usepackage{microtype}
\usepackage{booktabs}
\usepackage{array}
\usetikzlibrary{arrows.meta,positioning,fit,calc,shapes.geometric}

\graphicspath{{assets/}}

% ---------- Identidad visual UTP ----------
\definecolor{utpRojo}{HTML}{C42B35}
\definecolor{utpNegro}{HTML}{111111}
\definecolor{utpGris}{HTML}{5F5F5F}
\definecolor{utpGrisMedio}{HTML}{9A9A9A}
\definecolor{utpGrisClaro}{HTML}{F1F1F1}

\usetheme{default}
\setbeamercolor{normal text}{fg=utpNegro,bg=white}
\setbeamercolor{structure}{fg=utpRojo}
\setbeamercolor{frametitle}{fg=utpNegro,bg=white}
\setbeamercolor{title}{fg=utpNegro}
\setbeamercolor{subtitle}{fg=utpNegro}
\setbeamercolor{author}{fg=utpNegro}
\setbeamercolor{date}{fg=utpNegro}
\setbeamercolor{block title}{fg=white,bg=utpNegro}
\setbeamercolor{block body}{fg=utpNegro,bg=utpGrisClaro}
\setbeamercolor{alertblock title}{fg=white,bg=utpRojo}
\setbeamercolor{alertblock body}{fg=utpNegro,bg=utpRojo!7}
\setbeamercolor{exampleblock title}{fg=white,bg=utpGris}
\setbeamercolor{exampleblock body}{fg=utpNegro,bg=utpGris!8}
\setbeamercolor{terminal}{fg=white,bg=utpNegro}

\setbeamerfont{title}{series=\bfseries,size=\fontsize{19.5}{22.5}\selectfont}
\setbeamerfont{subtitle}{size=\fontsize{11.5}{13.5}\selectfont}
\setbeamerfont{date}{series=\bfseries,size=\fontsize{11.5}{13.5}\selectfont}
\setbeamerfont{author}{size=\fontsize{10.5}{12.5}\selectfont}
\setbeamerfont{frametitle}{series=\bfseries,size=\large}
\setbeamerfont{block title}{series=\bfseries,size=\small}
\setbeamerfont{block body}{size=\normalsize}

\setbeamertemplate{navigation symbols}{}
\setbeamertemplate{headline}{}
\setbeamertemplate{blocks}[default]
\setbeamertemplate{itemize item}{\color{utpRojo}\small\raise.15ex\hbox{$\blacktriangleright$}}
\setbeamertemplate{itemize subitem}{\color{utpGris}\scriptsize$\bullet$}
\setbeamertemplate{enumerate item}{\color{utpRojo}\insertenumlabel.}
\setbeamertemplate{background canvas}{\color{white}\rule{\paperwidth}{\paperheight}}

% Título de diapositiva: negrita negra a la izquierda, logo UTP a la derecha
\setbeamertemplate{frametitle}{%
  \begin{tikzpicture}[remember picture,overlay]
    \node[anchor=north east,inner sep=0pt]
      at ([xshift=-7mm,yshift=-5.5mm]current page.north east)
      {\includegraphics[width=17mm]{utp-logo.png}};
  \end{tikzpicture}%
  \vspace{1.2mm}%
  \begin{minipage}[t]{.80\textwidth}\raggedright\insertframetitle\end{minipage}\par
  \vspace{1.1mm}{\color{utpRojo}\rule{\textwidth}{1.1pt}}
}

% Solo numeración discreta; la plantilla UTP no lleva pie de página
\setbeamertemplate{footline}{%
  \hfill{\color{utpGrisMedio}\scriptsize\insertframenumber}\hspace{6mm}\vspace{2.5mm}%
}

% Portada UTP
\setbeamertemplate{title page}{%
  \begin{tikzpicture}[remember picture,overlay]
    \node[anchor=center,inner sep=0pt] at (current page.center) {%
      \includegraphics[width=\paperwidth,height=\paperheight]{utp-cover-bg.jpeg}%
    };
    \node[anchor=north,align=center,text=utpNegro,inner sep=0pt,
      text width=145mm,font=\usebeamerfont{title}]
      at ([yshift=-25.2mm]current page.north) {\inserttitle};
    \node[anchor=north,align=center,text=utpNegro,inner sep=0pt,
      text width=142mm,font=\usebeamerfont{subtitle}]
      at ([yshift=-38.2mm]current page.north) {\insertsubtitle};
    \node[anchor=north,align=center,text=utpNegro,inner sep=0pt,
      font=\usebeamerfont{date}]
      at ([yshift=-50.4mm]current page.north) {\insertdate};
    \node[anchor=north,align=center,text=utpNegro,inner sep=0pt,
      font=\usebeamerfont{author}]
      at ([yshift=-58.4mm]current page.north) {\insertauthor};
  \end{tikzpicture}%
}

% ---------- Comandos de apoyo didáctico ----------
\newcommand{\concepto}[2]{%
  \begin{block}{#1}\large #2\end{block}}
\newcommand{\pregunta}[1]{%
  \begin{alertblock}{Decide y explica}#1\end{alertblock}}
\newcommand{\analogia}[1]{%
  \begin{exampleblock}{En palabras simples}\small #1\end{exampleblock}}
\newcommand{\enlarealidad}[1]{%
  \begin{block}{Cómo se ve en la realidad}\small #1\end{block}}
\newcommand{\ojo}[1]{%
  \begin{alertblock}{Ojo con esto}\small #1\end{alertblock}}
\newcommand{\apoyo}{{\normalfont\footnotesize\color{utpGris}~· material de apoyo}}
\newcommand{\marca}[1]{\tikz[baseline=(m.base)]{%
  \node[rounded corners=1.5pt,fill=utpGrisClaro,draw=utpGrisMedio,inner xsep=3pt,inner ysep=1.5pt,
        font=\footnotesize] (m) {#1};}}
\newenvironment{terminal}{%
  \begin{beamercolorbox}[wd=\linewidth,sep=5pt,rounded=true]{terminal}%
  \ttfamily\small\ignorespaces}{%
  \end{beamercolorbox}}

\tikzset{
  caja/.style={draw=utpNegro, very thick, rounded corners=1.5pt, fill=white,
    align=center, minimum height=0.72cm, inner sep=5pt, font=\small},
  cajaazul/.style={caja, draw=utpGris, fill=white},
  cajaverde/.style={caja, draw=utpNegro, fill=utpGrisClaro},
  cajanaranja/.style={caja, draw=utpRojo!75, fill=utpRojo!7},
  flecha/.style={-{Stealth[length=2.3mm]}, very thick, draw=utpNegro},
  flechavuelta/.style={{Stealth[length=2.3mm]}-{Stealth[length=2.3mm]}, very thick, draw=utpNegro}
}

\graphicspath{{../../Sesion05/slides/assets/}}



\graphicspath{{assets/}{../../Sesion05/slides/assets/}}
\usepackage{colortbl}
\usepackage{listings}
% Un compose.yaml: las claves en rojo.
\lstset{basicstyle=\ttfamily\scriptsize,backgroundcolor=\color{utpGrisClaro},
  frame=none,aboveskip=4pt,belowskip=4pt,showstringspaces=false,
  columns=fullflexible,keepspaces=true,breaklines=true,
  keywords={services,build,image,ports,environment,volumes},
  keywordstyle=\color{utpRojo}\bfseries,
  commentstyle=\color{utpGris}\itshape,morecomment=[l]{\#},
  xleftmargin=5pt,xrightmargin=5pt,framesep=5pt}
% Una terminal: fondo oscuro, sin colorear palabras.
\lstdefinestyle{terminal}{basicstyle=\ttfamily\scriptsize\color{white},backgroundcolor=\color{utpNegro},
  keywords={},morecomment=[l]{\#},commentstyle=\color{utpGrisMedio},upquote=true}
\usepackage{upquote}
\usepackage{pifont}
\definecolor{okVerde}{HTML}{2E7D32}

\title{Integración de Servicios Cloud GCP}
\subtitle{Varios contenedores con Docker Compose}
\date{Semana 15 - Sesión 29}
\author{Rel Guzman}

% Íconos simples para los diagramas.
\tikzset{
 caja/.style={draw=utpRojo,thick,rounded corners=4pt,align=center},
 zona/.style={draw=utpGrisMedio,dashed,rounded corners=4pt},
 flecha/.style={-latex,thick},
}

\begin{document}

{\setbeamertemplate{footline}{}\maketitle}

% 2
\begin{frame}{Revisión: lo que vimos en la Sesión 28}
\begin{columns}[T,onlytextwidth]
\column{.46\textwidth}
\begin{block}{Dockerfile}
\small Los pasos para armar una imagen. \texttt{docker build -t cafe-api .} la construye.
\end{block}
\begin{block}{Puertos}
\small \texttt{-p 8080:8080} lleva el 8080 de tu computadora al 8080 del contenedor.
\end{block}
\begin{block}{Variables de entorno}
\small \texttt{-e PUERTO=9090} cambia la configuración sin cambiar la imagen.
\end{block}
\column{.50\textwidth}
\begin{alertblock}{Responde}
\small
1. ¿Qué instrucción del Dockerfile corre \textbf{al arrancar} el contenedor?\par\medskip
2. Ya tienes un contenedor en el 8080 de tu computadora. ¿Cómo corres otra copia de \texttt{cafe-api}?\par\medskip
3. Cambiaste \texttt{api.py}. ¿Qué hay que hacer para que el contenedor use el código nuevo?
\end{alertblock}
\end{columns}
\end{frame}

% 3
\begin{frame}{Revisión: la pregunta de hoy}
\begin{center}
\begin{tikzpicture}[font=\small]
  \node[caja,dashed,text width=2.2cm,minimum height=1.2cm] (nav) at (0,0) {\textbf{Navegador}\\\footnotesize\texttt{localhost:8080}};
  \node[caja,fill=utpGrisClaro,text width=2.2cm,minimum height=1.2cm] (web) at (3.8,0) {\textbf{web}\\\footnotesize la página};
  \node[caja,fill=utpRojo!30,text width=2.2cm,minimum height=1.2cm] (api) at (7.6,0) {\textbf{api}\\\footnotesize el menú y los pedidos};
  \node[caja,draw=okVerde,fill=okVerde!12,text width=2.2cm,minimum height=1.2cm] (db) at (11.4,0) {\textbf{db}\\\footnotesize la base de datos};
  \draw[flecha] (nav) -- (web);
  \draw[flecha,dashed] (web) -- node[above,font=\bfseries]{?} (api);
  \draw[flecha,dashed] (api) -- node[above,font=\bfseries]{?} (db);
\end{tikzpicture}
\end{center}
\vspace{3mm}
\begin{columns}[T,onlytextwidth]
\column{.32\textwidth}
\begin{block}{1. Compose}
\small Son tres contenedores. ¿Cómo se levantan todos con un comando?
\end{block}
\column{.32\textwidth}
\begin{block}{2. Redes}
\small ¿Cómo encuentra la API a la base de datos?
\end{block}
\column{.32\textwidth}
\begin{block}{3. Volúmenes}
\small Si se borra el contenedor de la base, ¿se pierden los pedidos?
\end{block}
\end{columns}
\end{frame}

% 4
\begin{frame}[c]
\begin{block}{Logro de la sesión}
Al finalizar, podrás explicar con tus palabras cómo \textbf{Docker Compose} levanta la página, la API y la base de Café Nube con un comando, cómo un contenedor encuentra a otro \textbf{por su nombre} y cómo un \textbf{volumen} guarda los pedidos aunque se borre el contenedor.
\end{block}
\end{frame}

% 5
\begin{frame}{Saberes previos}
\begin{itemize}\setlength\itemsep{14pt}
  \item En la Sesión 28, ¿qué comando escribiste para correr la API con su puerto publicado?
  \item Si la API de Café Nube necesita una base de datos, ¿cuántos \texttt{docker run} tendrías que escribir?
  \item ¿Qué pasa con un contenedor cuando corres \texttt{docker rm}?
\end{itemize}
\end{frame}

% 6
\begin{frame}[c]{Utilidad de esta sesión}
\begin{center}
\begin{minipage}{.84\textwidth}
{\large Una app real tiene varias piezas: Café Nube tiene una \textbf{página}, una \textbf{API} y una \textbf{base de datos}, cada una en su contenedor. Con un archivo, \texttt{compose.yaml}, las levantas todas con \textbf{un comando}, y los pedidos quedan guardados aunque apagues todo.}
\end{minipage}
\end{center}
\end{frame}

% 7
\begin{frame}[fragile]{Docker Compose: varias piezas, un comando}
\begin{columns}[T,onlytextwidth]
\column{.36\textwidth}
\begin{lstlisting}
services:
  web:
    build: ./web
    ports:
      - "8080:80"
  api:
    build: ./api
  db:
    image: postgres:17
\end{lstlisting}
{\scriptsize Resumido: el archivo completo está en la guía.}
\column{.60\textwidth}
\begin{block}{Servicio}
\small Cada pieza de la app: \texttt{web}, \texttt{api} y \texttt{db}. \texttt{build:} la construye con su Dockerfile; \texttt{image:} usa una imagen ya hecha.
\end{block}
\renewcommand{\arraystretch}{1.2}
\small
\begin{tabular}{@{}>{\ttfamily\footnotesize}l>{\raggedright\arraybackslash}p{4.0cm}@{}}
\toprule
\textsf{\textbf{Comando}} & \textbf{Qué hace}\\
\midrule
docker compose up -d & Levanta los tres servicios\\
docker compose ps & Lista sus contenedores\\
docker compose logs api & Los mensajes de \texttt{api}\\
docker compose down & Detiene y borra los contenedores\\
\bottomrule
\end{tabular}
\end{columns}
\end{frame}

% 8
\begin{frame}[fragile]{Redes: cada servicio, por su nombre}
\begin{center}
\begin{tikzpicture}[font=\small]
  \node[caja,dashed,text width=2.1cm,minimum height=1.0cm] (nav) at (0,0) {\textbf{Navegador}\\\footnotesize\texttt{localhost:8080}};
  \node[caja,fill=utpGrisClaro,text width=1.8cm,minimum height=1.0cm] (web) at (4.2,0) {\textbf{web}};
  \node[caja,fill=utpRojo!30,text width=1.8cm,minimum height=1.0cm] (api) at (7.9,0) {\textbf{api}};
  \node[caja,draw=okVerde,fill=okVerde!12,text width=1.8cm,minimum height=1.0cm] (db) at (11.6,0) {\textbf{db}};
  \node[zona,fit=(web)(db),inner sep=10pt,label={[font=\scriptsize,text=utpGris]below:red de Docker que crea Compose}] {};
  \draw[flecha] (nav) -- node[above,font=\scriptsize\ttfamily]{8080:80} (web);
  \draw[flecha] (web) -- node[above,font=\scriptsize\ttfamily]{api:8080} (api);
  \draw[flecha] (api) -- node[above,font=\scriptsize\ttfamily]{db:5432} (db);
\end{tikzpicture}
\end{center}
\vspace{2mm}
\begin{columns}[T,onlytextwidth]
\column{.48\textwidth}
\begin{block}{El nombre del servicio es su dirección}
\small Compose pone los tres contenedores en una \textbf{red}. Ahí, la API está en \texttt{api} y la base en \texttt{db}. Solo \texttt{web} publica un puerto, para el navegador.
\end{block}
\column{.48\textwidth}
\begin{alertblock}{Ojo}
\small Dentro de un contenedor, \texttt{localhost} es \textbf{ese mismo} contenedor. Por eso la API no busca la base en \texttt{localhost}, sino en \texttt{db}.
\end{alertblock}
\end{columns}
\end{frame}

% 9
\begin{frame}[fragile]{Volúmenes: los pedidos sobreviven}
\begin{block}{Volumen}
\small Un espacio de guardado aparte del contenedor. Lo que la base escribe ahí \textbf{queda aunque se borre el contenedor}.
\end{block}
\begin{columns}[T,onlytextwidth]
\column{.48\textwidth}
\begin{lstlisting}
  db:
    image: postgres:17
    environment:
      POSTGRES_PASSWORD: cafe
    volumes:
      - datos:/var/lib/postgresql/data

volumes:
  datos:
\end{lstlisting}
{\scriptsize\texttt{environment:} son las variables, lo que antes era \texttt{-e}.}
\column{.48\textwidth}
\renewcommand{\arraystretch}{1.3}
\small
\begin{tabular}{@{}>{\ttfamily\footnotesize}l>{\raggedright\arraybackslash}p{3.0cm}@{}}
\toprule
\textsf{\textbf{Comando}} & \textbf{Los pedidos}\\
\midrule
docker compose down & \textcolor{okVerde}{\textbf{Siguen}}: el volumen queda\\
docker compose down -v & \textcolor{utpRojo}{\textbf{Se borran}}: borra el volumen\\
\bottomrule
\end{tabular}
\end{columns}
\end{frame}

% 10
\begin{frame}{Cierre: preguntas para comprobar}
\begin{block}{Responde con tus palabras}
\small
1. ¿Qué comando levanta todos los servicios de \texttt{compose.yaml}?\\\medskip
2. Desde el contenedor \texttt{api}, ¿qué dirección usas para llegar a la base: \texttt{localhost:5432} o \texttt{db:5432}?\\\medskip
3. Corres \texttt{docker compose down} y luego \texttt{up -d}. ¿Siguen los pedidos? ¿Por qué?
\end{block}
\end{frame}

% 11
\begin{frame}{Referencias}
\scriptsize
\begin{itemize}\setlength\itemsep{4pt}
  \item Ruta \textit{Despliegue con Docker}, curso 2, \textit{Varios contenedores con Compose}: varias piezas, \texttt{compose.yaml}, redes entre contenedores, volúmenes y el escenario Café Nube.
  \item Docker Compose: \texttt{docs.docker.com}. Imágenes \texttt{nginx} y \texttt{postgres}: \texttt{hub.docker.com}.
  \item Material previo del curso: Sesión 28 (Dockerfile, puertos y variables de entorno).
  \item Las respuestas de las preguntas están en \texttt{../respuestas.md}.
\end{itemize}
\end{frame}

\end{document}
